% Einheit 4 von 5 – Grundwert berechnen (bank/prozentrechnung/e4.jsonl, zusammenbau v0.1)
%% TODO Zweigzeile Teil 1: was hier gelernt wird (Satz in Schülersprache aus dem Einheitstitel) – Einheit 4
\einheitenkopf[e4]{Einheit 4 von 5 · Grundwert berechnen}
\zweigzeile{ab Kl. 7 · P10 oft}
\setcounter{aufgabe}{19}

%% TODO Titel: Ich-kann-Satz fehlt in der Bank (Platzhalter: Kettenname) – Nr. 20
%% TODO Anweisung: ein Satz über der ersten Teilaufgabe fehlt in der Bank – Nr. 20
\begin{aufgabe}{Was ist gegeben, was gesucht?}
\begin{teile}
\teil $30\,\%$ der Klasse sind $9$ Kinder. Wie viele Kinder hat die Klasse? – Was ist gesucht?\\ \kreuz{Prozentwert}\\ \kreuz{Prozentsatz}\\ \kreuz{Grundwert}
\end{teile}
\end{aufgabe}

%% TODO Titel: Ich-kann-Satz fehlt in der Bank (Platzhalter: Kettenname) – Nr. 21
%% TODO Anweisung: ein Satz über der ersten Teilaufgabe fehlt in der Bank – Nr. 21
\begin{aufgabe}{Grundwert}
%% TODO \rechenplatz{n}: Schrittzahl des Grundfalls fehlt in der Bank (nur bei fester Schreibform, 2.3 b)
\begin{teile}
\teil Grau sind $30\,\%$, das sind $6$ Kästchen – wie viele Kästchen hat der ganze Streifen?

\streifen[0]{30}{$0\,\%$}{$100\,\%$}
\leerfeld[Kästchen]
\teil $50\,\%$ sind $35$ kg – wie viel ist das Ganze? \leerfeld[kg]
\teil $3\,\%$ sind $12$ kg – wie viel ist das Ganze? \leerfeld[kg]
\teil $36\,\%$ sind $81$ kg – wie viel ist das Ganze? \leerfeld[kg]
\teil $21$ Kinder fahren mit dem Rad, das sind $75\,\%$ der Klasse – wie viele Kinder hat die Klasse? \leerfeld
\teil $40\,\%$ von $65$ Kindern gehen in die AG – wie viele? \leerfeld
\teil Beim Kauf einer Mütze spart Mira $4\,€$, das sind $25\,\%$ des alten Preises. Wie viel kostete die Mütze vorher? \hfill (P10 2025 OS) \\ \leerfeld[€]
\end{teile}
\end{aufgabe}

%% TODO Titel: Ich-kann-Satz fehlt in der Bank (Platzhalter: Kettenname) – Nr. 22
%% TODO Anweisung: ein Satz über der ersten Teilaufgabe fehlt in der Bank – Nr. 22
\begin{aufgabe}{Grundwert – Fehler finden, Begründen, Darstellungswechsel, Anwendung}
\begin{teile}
\teil „$30\,\%$ sind $12$ kg – wie viel ist das Ganze?“ Sven: \rechnung{0{,}3 \cdot 12 &= 3{,}6\text{ kg}} Fehler? Wie heißt es richtig?
\teil Warum rechnet man beim Grundwert am Ende mal $100$?
\teil Grau sind $40\,\%$, das sind $18$ kg. Trage die Werte an den Streifen, lies das Ganze ab.

\streifen{40}{$0$ kg}{?}
\leerfeld[kg]
\teil Tims Handy zeigt $35\,\%$ Akku; laut Anzeige reicht das noch für $7$ Stunden Musik. Wie lange hält ein voller Akku? Reicht er für eine Busfahrt von $18$ Stunden?
\end{teile}
\end{aufgabe}
